\documentclass[letterpaper,11pt]{article}

\usepackage{latexsym}
\usepackage[empty]{fullpage}
\usepackage{titlesec}
\usepackage{marvosym}
\usepackage[usenames,dvipsnames]{color}
\usepackage{verbatim}
\usepackage{enumitem}
\usepackage[hidelinks]{hyperref}
\usepackage{fancyhdr}
\usepackage[english]{babel}
\usepackage[utf8]{inputenc}
\usepackage{tabularx}
\input{glyphtounicode}

\pagestyle{fancy}
\fancyhf{}
\fancyfoot{}
\renewcommand{\headrulewidth}{0pt}
\renewcommand{\footrulewidth}{0pt}

\addtolength{\oddsidemargin}{-0.5in}
\addtolength{\evensidemargin}{-0.5in}
\addtolength{\textwidth}{1in}
\addtolength{\topmargin}{-0.88in}
\addtolength{\textheight}{1.9in}

\urlstyle{same}
\raggedbottom
\raggedright
\setlength{\tabcolsep}{0in}

\titleformat{\section}{
  \vspace{-4pt}\scshape\raggedright\large
}{}{0em}{}[\color{black}\titlerule \vspace{-5pt}]

\pdfgentounicode=1
\newcommand{\resumeName}[1]{\textbf{\Huge \scshape #1} \\ \vspace{1pt}}

\newcommand{\resumeItem}[1]{
  \item\small{
    {#1 \vspace{-2pt}}
  }
}

\newcommand{\resumeSubheading}[4]{
  \vspace{-2pt}\item
    \begin{tabular*}{0.97\textwidth}[t]{l@{\extracolsep{\fill}}r}
      \textbf{#1} & #2 \\
      \textit{\small#3} & \textit{\small #4} \\
    \end{tabular*}\vspace{-7pt}
}

\newcommand{\resumeSubSubheading}[2]{
    \item
    \begin{tabular*}{0.97\textwidth}{l@{\extracolsep{\fill}}r}
      \textit{\small#1} & \textit{\small #2} \\
    \end{tabular*}\vspace{-7pt}
}

\newcommand{\resumeProjectHeading}[2]{
    \item
    \begin{tabular*}{0.97\textwidth}{l@{\extracolsep{\fill}}r}
      \small#1 & #2 \\
    \end{tabular*}\vspace{-7pt}
}

\newcommand{\resumePatentHeading}[3]{
    \item
    \begin{tabularx}{0.97\textwidth}[t]{@{}>{\raggedright\arraybackslash}X r@{}}
      \small\textbf{#1} -- \textit{#2} & \small #3 \\
    \end{tabularx}\vspace{-7pt}
}

\newcommand{\resumeSubItem}[1]{\resumeItem{#1}\vspace{-4pt}}

\renewcommand\labelitemii{$\vcenter{\hbox{\tiny$\bullet$}}$}

\newcommand{\resumeSubHeadingListStart}{\begin{itemize}[leftmargin=0.15in, label={}]}
\newcommand{\resumeSubHeadingListEnd}{\end{itemize}}
\newcommand{\resumeItemListStart}{\begin{itemize}[itemsep=-1pt,topsep=2pt]}
\newcommand{\resumeItemListEnd}{\end{itemize}\vspace{-5pt}}

\begin{document}

\begin{center}
  \resumeName{Harinandan Kotamsetti}
    \small Cupertino, CA $|$ +1 (669) 264-9458 $|$
    \href{mailto:hareee234@gmail.com}{\underline{hareee234@gmail.com}} $|$
    \href{https://www.linkedin.com/in/hareee234/}{\underline{linkedin.com/in/hareee234}} $|$
    \href{https://github.com/AwesomenessReborn}{\underline{github.com/AwesomenessReborn}}
\end{center}

\section{Education}
  \resumeSubHeadingListStart
    \resumeSubheading
      {\textbf{San Jose State University}}{San Jose, CA}
      {Bachelor of Science in Computer Engineering}{Expected Dec 2026}
      \resumeItemListStart
        \resumeItem{\textbf{Coursework:} OS, Computer Architecture, C/C++, FPGA, Verilog, Linear Algebra, x86 \& MIPS Assembly}
      \resumeItemListEnd
  \resumeSubHeadingListEnd

\section{Experience}
  \resumeSubHeadingListStart
    \resumeSubheading
      {\textbf{STMicroelectronics}}{Santa Clara, CA}
      {\textbf{Software \& Algorithm Intern}}{Mar 2025 -- Present}
      \resumeItemListStart
        \resumeItem{
          \textbf{Quantized a 103K-parameter walking-speed CNN-LSTM to INT8 and deployed on Snapdragon NPU via Qualcomm QNN} (0.059 m/s MAE): built ring-buffered \textbf{C++ IMU preprocessing at 208 Hz} and an \textbf{Android JNI} bridge, running at \textbf{$\sim$15 ms per 2.5 s window}; demoed at CES 2026.
        }
        \resumeItem{
          Developed a \textbf{C/Python regression harness} validating exported-model outputs against \textbf{Python references across TFLite, ONNX, and QNN} targets, enforcing numerical correctness across model updates.
        }
        \resumeItem{
          Ported an \textbf{887K-parameter VR head-motion model} from \textbf{PyTorch $\rightarrow$ ONNX $\rightarrow$ QAIRT $\rightarrow$ Hexagon NPU}, matching PyTorch output within \textbf{0.056$^\circ$ on-device} after fixing converter and tensor-layout failures.
        }
        \resumeItem{
          \textbf{Deployed always-on fall detection and wrist-tilt recognition entirely in-sensor} (ST ISPU/MLC, \textbf{zero host-MCU compute}); sustained \textbf{104 Hz sensor fusion at 3.08 ms worst case} (32\% of budget) in \textbf{$\sim$20 KB code / 1.1 KB RAM}.
        }
        \resumeItem{
          \textbf{Owned the head-motion training pipeline for 45+ GPU-hour runs}, increasing throughput \textbf{2.5$\times$} with bf16 + \texttt{torch.compile} and adding per-horizon best-checkpoint saving after finding $\sim$79\% of evaluated checkpoints were not retained.
        }
      \resumeItemListEnd

    \resumeSubheading
      {\textbf{Spartan Racing (Formula SAE)}}{San Jose, CA}
      {\textbf{Embedded Systems \& Software Developer}}{Aug 2023 -- Feb 2025}
      \resumeItemListStart
        \resumeItem{
          \textbf{Built a Python hardware-in-the-loop test suite} on a breakout board emulating an accumulator module, validating OV/UV thresholds and fault logic without a live 400 V pack.
        }
        \resumeItem{
          \textbf{Cut STM32 BMS polling latency 14$\times$ (7.0 s $\rightarrow$ 0.5 s)} by eliminating blocking CAN/thermistor delays and adding mailbox-aware CAN transmit logic.
        }
        \resumeItem{
          \textbf{Implemented state-of-charge estimation} (Coulomb counting with open-circuit-voltage lookup-table correction) for a \textbf{400 V, 768-cell BMS on STM32}, reaching \textbf{3\% MAE}.
        }
      \resumeItemListEnd
  \resumeSubHeadingListEnd

\section{Patents \& Invention Disclosures}
  \resumeSubHeadingListStart
    \resumePatentHeading
      {Glasses with Wellness and Health Monitoring}
      {STMicroelectronics, Co-inventor}
      {Filed May 2026}
    \resumePatentHeading
      {Motion-to-Photon Latency Compensation for XR/AR/VR Glasses}
      {STMicroelectronics, Co-inventor}
      {Invention Disclosure, 2026}
  \resumeSubHeadingListEnd

\section{Projects}
  \resumeSubHeadingListStart
    \resumeProjectHeading
      {\textbf{MERIT UAV --- Autonomous Surveillance Drone} $|$
        \emph{Python, YOLOv8n, CUDA, Jetson, MAVLink}}{2026}
      \resumeItemListStart
        \resumeItem{
          \textbf{Implemented and flight-tested autonomous MAVLink control} (takeoff, altitude, position, yaw) from a Jetson Orin Nano to a Pixhawk.
        }
        \resumeItem{
          Ran \textbf{4 Hz GPS telemetry alongside 30 FPS CUDA YOLOv8n detection}, geotagging detections and gating uploads on GPS lock; vision pipeline ground-tested.
        }
      \resumeItemListEnd

    \resumeProjectHeading
      {\textbf{Nemotron-H Fine-Tuning --- Kaggle NLP Competition} $|$
        \emph{Python, PyTorch, LoRA, llama.cpp}}{2026}
      \resumeItemListStart
        \resumeItem{
          \textbf{Improved controlled task accuracy by +19.6--22.5 absolute percentage points} with a LoRA fine-tuning/evaluation pipeline for Nemotron-H 4B, training \textbf{20.2M / 3.99B parameters (0.5\%)} on \textbf{8,000+ self-generated reasoning traces} with assistant-only label masking.
        }
      \resumeItemListEnd
  \resumeSubHeadingListEnd

\section{Technical Skills}
  \resumeSubHeadingListStart
    \small{\item{
      \textbf{Languages}{: Python, C++, C, Kotlin, Java, SQL} \\
      \textbf{ML \& Deployment}{: PyTorch, TFLite, ONNX, Qualcomm QNN/QAIRT, TensorFlow, INT8 quantization, LoRA} \\
      \textbf{Mobile \& Systems}{: Android NDK/JNI, Gradle, adb/logcat, CMake, STM32, CAN, Git, GitHub Actions} \\
      \textbf{Sensors \& Real-Time ML}{: IMU, sensor fusion, on-device ML, real-time inference, Butterworth filtering, FFT} \\
      \textbf{Validation \& Analysis}{: C/Python regression testing, HIL testing, numerical validation, NumPy, Pandas, SciPy}
    }}
  \resumeSubHeadingListEnd

\end{document}
